\documentclass[Riemann_Hypothesis.tex]{subfiles}
\begin{document}
\pagebreak
\color{red}
\part{Résolution de la conjecture}
\color{black}

\color{red}
\section{Premières définitions et espaces de Weil}
\color{black}

\begin{DefNef}[label=df:EspaceSchwartz_fr]{Espace de Schwartz}
    Soit $\mathbb{R}$ le corps des nombres réels. \par
    D'après les travaux de l'analyse fonctionnelle introduits par Laurent 
    Schwartz~\cite{schwartz1966_distributions}, on appelle espace de Schwartz, 
    que l'on notera $\mathcal{S}(\mathbb{R})$, l'espace vectoriel des 
    fonctions $f \in C^\infty(\mathbb{R}, \mathbb{C})$ dites à décroissance 
    rapide. \par
    On dira qu'une fonction $f$ appartient à l'espace $\mathcal{S}(\mathbb{R})$ si et seulement si, pour tout couple d'entiers $(\alpha, \beta) \in \mathbb{N}^2$, on a :
    \begin{equation*}
        \sup_{x \in \mathbb{R}} \left| x^\alpha f^{(\beta)}(x) \right| \;<\; +\infty
    \end{equation*}
    où $f^{(\beta)}$ désigne la dérivée d'ordre $\beta$ de la fonction $f$.
\end{DefNef}

\begin{DefNef}[label=df:EspaceWeil_fr]{Espace des fonctions tests de Weil}
    Soit $\mathbb{R}_+^*$ le groupe multiplicatif des réels strictement positifs. \par
    D'après les travaux d'André Weil~\cite{weil1952_formules}, 
    on appelle espace des fonctions tests de Weil, que l'on notera $\mathcal{W}$, l'ensemble des fonctions $f \in C^\infty(\mathbb{R}_+^*, \mathbb{C})$ telles que la fonction $g : t \mapsto f(e^t)$ appartient à l'espace de Schwartz $\mathcal{S}(\mathbb{R})$. \par
    Afin de formuler une condition équivalente, on introduit l'opérateur différentiel $D$ défini sur $C^\infty(\mathbb{R}_+^*, \mathbb{C})$ par :
    \begin{equation*}
        D(f)(x) = x f'(x)
    \end{equation*}
    On dira de manière équivalente que $f \in \mathcal{W}$ si et seulement si, pour tout couple d'entiers $(k, m) \in \mathbb{N}^2$, on a :
    \begin{equation*}
        \sup_{x \in \mathbb{R}_+^*} \left| (\ln x)^k D^m(f)(x) \right| \;<\; +\infty
    \end{equation*}
    où $D^m$ désigne la $m$-ième itérée de l'opérateur $D$.
\end{DefNef}


\begin{DefNef}[label=df:MesureHaar_fr]{Mesure de Haar multiplicative}
    D'après les travaux d'Alfred Haar (voir Bourbaki~\cite{Haar_bourbaki_integration_chap7}), 
    sur tout groupe topologique localement compact, 
    il existe une mesure de Radon régulière et non triviale, unique à une constante
    multiplicative près, qui est invariante par translation. \par
    Soit $G = \left(\mathbb{R}_+^*, \times\right)$ le groupe topologique multiplicatif 
    des réels strictement positifs. \par
    On appelle \textbf{mesure de Haar} du groupe $G$, que l'on notera $d^\times x$, 
    l'unique mesure (à un facteur près) invariante par l'action de multiplication. 
    Elle est définie par :
    \begin{equation*}
        d^\times x \;=\; \frac{\mathrm{d}x}{x}
    \end{equation*}
    où $\mathrm{d}x$ désigne la mesure de Lebesgue usuelle sur $\mathbb{R}$. \par
    En effet, soit $a$ un réel strictement positif. \par
    On pose l'application de translation à gauche sur le groupe, définie par la 
    multiplication $x \mapsto ax$. \par
    Par différenciation de cette application, on a $\mathrm{d}(ax) = a \, \mathrm{d}x$. \par
    On évalue la mesure sur l'élément translaté :
    \begin{equation*}
        d^\times (ax) \;=\; \frac{\mathrm{d}(ax)}{ax} \;=\; \frac{a \, \mathrm{d}x}{ax} \;=\; \frac{\mathrm{d}x}{x}
    \end{equation*}
    Comme on a $d^\times (ax) = d^\times x$, on en déduit que la mesure ainsi
    construite est bien invariante par dilatation.

    C'est cette invariance qui assure la compatibilité structurelle avec la
    transformée de Mellin et la forme hermitienne de Weil.
\end{DefNef}


\begin{lemma}[label=lm:DensiteWeil_fr]{Densité de l'espace de Weil}
    Soit $L^2(\mathbb{R}_+^*, d^\times x)$ l'espace de Hilbert des fonctions à
    valeurs complexes de carré intégrable pour la mesure de Haar. \par
    L'espace des fonctions tests de Weil $\mathcal{W}$ est dense dans
    $L^2(\mathbb{R}_+^*, d^\times x)$.
\end{lemma}

\begin{proof}
    \begin{part_proof}
        Soit $\Phi : \mathbb{R} \to \mathbb{R}_+^*$ l'application définie par $\Phi(t) = e^t$. \par
        Comme la fonction exponentielle est un difféomorphisme de $\mathbb{R}$ sur $\mathbb{R}_+^*$, on pose l'opérateur de composition $J$, qui à toute fonction $f \in L^2(\mathbb{R}_+^*, d^\times x)$ associe la fonction $J(f) = f \circ \Phi$. \par
        On évalue la norme au carré de $J(f)$ dans l'espace usuel $L^2(\mathbb{R}, dt)$ :
        \begin{equation*}
            \int_{\mathbb{R}} \left| J(f)(t) \right|^2 \, dt \;=\; \int_{\mathbb{R}} \left| f(e^t) \right|^2 \, dt
        \end{equation*}
        On effectue le changement de variable $x = e^t$. \par
        Comme on a $dx = e^t \, dt = x \, dt$, on en déduit que $dt = \frac{dx}{x} = d^\times x$. \par
        En appliquant ce changement de variable aux bornes et à la mesure d'intégration, on obtient :
        \begin{equation*}
            \int_{\mathbb{R}} \left| f(e^t) \right|^2 \, dt \;=\; \int_{0}^{+\infty} \left| f(x) \right|^2 \, \frac{dx}{x} \;=\; \int_{\mathbb{R}_+^*} \left| f(x) \right|^2 \, d^\times x
        \end{equation*}
        Comme les normes sont conservées, on en déduit que l'opérateur $J$
        est une isométrie surjective de l'espace de Hilbert
        $L^2(\mathbb{R}_+^*, d^\times x)$ vers l'espace de Hilbert
        $L^2(\mathbb{R}, dt)$. \medskip\par
        Par définition de l'espace $\mathcal{W}$
        (voir Définition \ref{df:EspaceWeil_fr}),
        l'image de $\mathcal{W}$ par l'isomorphisme $J$ est exactement
        l'espace de Schwartz $\mathcal{S}(\mathbb{R})$. \par
        D'après les théorèmes fondamentaux de l'analyse fonctionnelle
        (voir les travaux de Laurent Schwartz~\cite{schwartz1966_distributions}), 
        l'espace $\mathcal{S}(\mathbb{R})$ est dense dans $L^2(\mathbb{R}, dt)$. \par
        Comme on a la densité de $J(\mathcal{W})$ dans $L^2(\mathbb{R}, dt)$ et que $J^{-1}$ 
        est une application continue (puisque $J$ est une isométrie), on en déduit que 
        $\mathcal{W} = J^{-1}(\mathcal{S}(\mathbb{R}))$ est dense dans 
        $L^2(\mathbb{R}_+^*, d^\times x)$.
    \end{part_proof}
\end{proof}

\begin{DefNef}[label=df:TransformeeMellin_fr]{Transformée de Mellin}
    Soit $f$ une fonction à valeurs complexes, définie et localement intégrable sur le groupe multiplicatif $\mathbb{R}_+^*$. \par
    D'après les travaux de Robert Hjalmar 
    {Mellin}, repris par Jean Dieudonné (voir~\cite{Mellin_dieudonne_calcul_infinitesimal}), on appelle transformée de 
    Mellin de $f$, 
    que l'on notera $\mathcal{M}[f]$, la fonction de la variable complexe $s$ définie par :
    \begin{equation*}
        \mathcal{M}[f](s) \;=\; \int_{0}^{+\infty} f(x) x^{s-1} \, dx \;=\; \int_{0}^{+\infty} f(x) x^s \, \frac{dx}{x}
    \end{equation*}
    pour tout $s \in \mathbb{C}$ tel que cette intégrale soit absolument convergente. \par
    On pose $\mathcal{D}_f \subset \mathbb{C}$ le domaine de convergence absolue
    de cette intégrale, qui constitue usuellement une bande verticale dans le
    plan complexe.
\end{DefNef}

\begin{proposition}[label=pr:MellinHolomorpheWeil_fr]{Holomorphie de la transformée de Mellin sur l'espace de Weil}
    Pour toute fonction $f \in \mathcal{W}$, sa transformée de Mellin
    $\mathcal{M}[f]$ est une fonction entière, donc
    holomorphe sur l'ensemble du plan complexe $\mathbb{C}$.
\end{proposition}

\begin{proof}
    \begin{part_proof}
        Soit $f \in \mathcal{W}$. \par
        Soit $s \in \mathbb{C}$. \par
        On pose $\sigma = \Re(s)$ la partie réelle du nombre complexe $s$. \par
        D'après la définition de l'espace de Weil (voir Définition \ref{df:EspaceWeil_fr}), 
        on sait que la fonction $g : t \mapsto f(e^t)$ appartient à l'espace de Schwartz 
        $\mathcal{S}(\mathbb{R})$. \par
        On en déduit que pour tout entier $N \in \mathbb{N}$, il existe une constante 
        $C_N > 0$ telle que pour tout $t \in \mathbb{R}$, on a :
        \begin{equation*}
            |g(t)| \;=\; |f(e^t)| \;\leq\; \frac{C_N}{1 + e^{N|t|}}
        \end{equation*}
        En effectuant le changement de variable $x = e^t$, comme on a $x \in \mathbb{R}_+^*$, 
        on en déduit que pour tout entier $N \in \mathbb{N}$, il existe un réel $M_N > 0$ 
        tel que :
        \begin{equation*}
            |f(x)| \;\leq\; M_N \min\left(x^N, x^{-N}\right)
        \end{equation*}
        Cette majoration prouve que la fonction $f$ possède une décroissance plus rapide que 
        toute puissance de $x$ lorsque $x \to 0$, et plus rapide que toute puissance de 
        l'inverse de $x$ lorsque $x \to +\infty$. \par
        On évalue le module de l'intégrande définissant la transformée de Mellin :
        \begin{equation*}
            \left| f(x) x^{s-1} \right| \;=\; |f(x)| x^{\sigma-1}
        \end{equation*}
        Comme on a la majoration précédente pour tout entier $N$, on peut choisir
        $N > |\sigma - 1| + 1$, et on en déduit que la fonction
        $x \mapsto |f(x)| x^{\sigma-1}$ est intégrable sur $\mathbb{R}_+^*$. \medskip\par

        Ainsi, pour tout $s \in \mathbb{C}$, l'intégrale définissant
        $\mathcal{M}[f](s)$ est absolument convergente. \par
        De plus, la fonction $s \mapsto f(x) x^{s-1}$ est entière
        pour tout $x \in \mathbb{R}_+^*$. \par
        D'après le théorème d'holomorphie des intégrales à paramètres,
        reposant sur le théorème de convergence dominée d'Henri 
        Lebesgue~\cite{lebesgue1902_these}, on peut dominer localement l'intégrande
        indépendamment de $s$ sur tout compact de $\mathbb{C}$.
    \end{part_proof}

    \begin{part_proof}
        D'après le théorème d'holomorphie des intégrales à paramètres, reposant 
        sur la théorie de la mesure 
        d'Henri Lebesgue~\cite{lebesgue1902_these}, pour qu'une fonction définie par une 
        intégrale soit holomorphe sur un ouvert $U$ de $\mathbb{C}$, il suffit que 
        l'intégrande soit mesurable par rapport à la variable d'intégration, holomorphe 
        par rapport au paramètre sur $U$, et qu'il admette une domination par une fonction 
        intégrable indépendante du paramètre sur tout compact inclus dans $U$. \par
        Soit $s_0 \in \mathbb{C}$. \par
        On pose $K$ un voisinage compact de $s_0$ dans le plan complexe $\mathbb{C}$, 
        tel qu'un disque fermé centré en $s_0$ de rayon strictement positif. \par
        Comme on a établi précédemment l'existence d'une fonction $\varphi$ intégrable sur 
        $\mathbb{R}_+^*$ telle que pour tout $s \in K$ et pour tout $x \in \mathbb{R}_+^*$, 
        on a $\left| f(x) x^{s-1} \right| \leq \varphi(x)$, on en déduit que les hypothèses 
        du théorème d'holomorphie sont vérifiées sur l'intérieur de ce compact $K$. \par
        On en déduit que la fonction $\mathcal{M}[f]$ est holomorphe sur l'intérieur de $K$. \par
        Comme le point $s_0$ appartient à l'intérieur de $K$, on en déduit que la fonction 
        $\mathcal{M}[f]$ est complexe-dérivable au point $s_0$. \par
        Comme on a ce résultat pour tout point $s_0 \in \mathbb{C}$, on en déduit que la 
        fonction $\mathcal{M}[f]$ est complexe-dérivable en tout point du plan complexe. \par
        Par définition, on en déduit que la fonction $\mathcal{M}[f]$ est holomorphe sur 
        l'ensemble du plan complexe $\mathbb{C}$, c'est-à-dire qu'elle est une fonction 
        entière.
    \end{part_proof}
\end{proof}

\color{red}
\section{Forme hermitienne de Weil et critère de positivité de Weil}
\color{black}

\begin{DefNef}[label=df:FormeHermitienneWeil_fr]{Forme hermitienne de Weil}
    Soit $\mathcal{W}$ l'espace des fonctions tests de Weil sur le groupe 
    multiplicatif $\left(\mathbb{R}_+^*, \times\right)$. \par
    Soient $f$ et $g$ deux fonctions de $\mathcal{W}$. \par
    On pose $g^*$ la fonction involutive définie pour tout $x \in \mathbb{R}_+^*$ par :
    \begin{equation*}
        g^*(x) \;=\; \frac{1}{x} \, \overline{g\left(\frac{1}{x}\right)}
    \end{equation*}
    On note $\Phi = f * g^*$ le produit de convolution multiplicatif de $f$ et $g^*$. \par
    D'après les formules explicites d'André Weil, on appelle 
    \textbf{forme hermitienne de Weil}, que l'on notera $\langle f, g \rangle_{\mathcal{HP}}$, la forme sesquilinéaire définie par l'action de la distribution géométrique sur la fonction test $\Phi$ :
    \begin{equation*}
        \langle f, g \rangle_{\mathcal{HP}} \;:=\; W_{\text{pol}}(\Phi) - W_{\text{arith}}(\Phi) - W_{\text{arch}}(\Phi)
    \end{equation*}
    où les différentes fonctionnelles agissant sur $\Phi$ sont explicitées analytiquement par :
    \begin{align*}
        W_{\text{pol}}(\Phi) \;&=\; \mathcal{M}[\Phi](1) + \mathcal{M}[\Phi](0) \\
        W_{\text{arith}}(\Phi) \;&=\; \sum_{p \in \mathbb{P}} \sum_{n=1}^{\infty} \frac{\ln(p)}{p^{n/2}} \left( \Phi(p^n) + \Phi\left(\frac{1}{p^n}\right) \right) \\
        W_{\text{arch}}(\Phi) \;&=\; \left( \ln(\pi) + \gamma \right) \Phi(1) + \int_{1}^{\infty} \frac{\Phi(x) + \Phi\left(\frac{1}{x}\right) - 2\Phi(1)}{x - \frac{1}{x}} \, \frac{\mathrm{d}x}{x}
    \end{align*}
    avec $\mathcal{M}$ désignant la transformée de Mellin, $\mathbb{P}$ l'ensemble des nombres premiers et $\gamma$ la constante d'Euler-Mascheroni.
\end{DefNef}

\begin{proposition}[label=pr:CriterePositiviteWeil_fr]{Critère de positivité de Weil}
    Soit $\mathcal{Z}$ l'ensemble des zéros non triviaux de
    la fonction Zeta de Riemann.\par
    D'après le théorème d'André Weil~\cite{weil1952_formules}, 
    la forme hermitienne $\langle \cdot, \cdot \rangle_{\mathcal{HP}}$ 
    possède les propriétés suivantes : \par
    Pour toute fonction $f \in \mathcal{W}$, l'évaluation de cette forme sur la
    diagonale s'exprime analytiquement comme une somme sur l'ensemble
    $\mathcal{Z}$:
    \begin{equation*}
        \langle f, f \rangle_{\mathcal{HP}} \;=\; \sum_{\rho \in \mathcal{Z}} \mathcal{M}[f](\rho) \, \overline{\mathcal{M}[f](1 - \overline{\rho})}
    \end{equation*}
    où $\mathcal{M}[f]$ désigne la transformée de Mellin de la fonction test $f$. \par
    De plus, l'Hypothèse de Riemann est vérifiée si et seulement si, pour toute fonction 
    non nulle $f \in \mathcal{W}$, on a :
    \begin{equation*}
        \langle f, f \rangle_{\mathcal{HP}} \;\geq\; 0
    \end{equation*}
\end{proposition}

\begin{proof}
    \begin{part_proof}
        Démontrons en premier lieu l'égalité donné dans la proposition. \par
        Soit $f \in \mathcal{W}$. \par
        On pose $\Phi = f * f^*$. Par définition de la forme hermitienne de Weil,
        on a:
        \begin{equation*}
            \langle f, f \rangle_{\mathcal{HP}} \;=\; W_{\text{pol}}(\Phi) - W_{\text{arith}}(\Phi) - W_{\text{arch}}(\Phi)
        \end{equation*}
        D'après les formules explicites d'André Weil, la somme des contributions
        polaire, arithmétique et archimédienne évaluée sur une fonction test
        équivaut à la somme de sa transformée de Mellin évaluée sur l'ensemble
        $\mathcal{Z}$.
        On en déduit que:
        \begin{equation*}
            \langle f, f \rangle_{\mathcal{HP}} \;=\; \sum_{\rho \in \mathcal{Z}}
            \mathcal{M}[\Phi](\rho)
        \end{equation*}
        Par propriété du produit de convolution multiplicatif associé à la mesure
        de Haar $\frac{\mathrm{d}x}{x}$, la transformée de Mellin transforme la
        convolution en produit usuel. Pour tout nombre complexe $s$, on a:
        \begin{equation*}
            \mathcal{M}[\Phi](s) \;=\; \mathcal{M}[f](s) \mathcal{M}[f^*](s)
        \end{equation*}
        Exprimons la transformée de Mellin de l'involution $f^*$. On a:
        \begin{equation*}
            \mathcal{M}[f^*](s) \;=\; \int_{0}^{\infty} f^*(x) x^{s-1} \,
            \mathrm{d}x \;=\; \int_{0}^{\infty} \frac{1}{x}
            \overline{f\left(\frac{1}{x}\right)} x^{s-1} \, \mathrm{d}x
        \end{equation*}
        On effectue le changement de variable $y = \frac{1}{x}$. On en déduit
        que $\mathrm{d}x = - \frac{1}{y^2} \mathrm{d}y$. Comme les bornes
        d'intégration s'inversent, on a:
        \begin{equation*}
            \mathcal{M}[f^*](s) \;=\; \int_{0}^{\infty} y \overline{f(y)} y^{1-s}
            \frac{1}{y^2} \, \mathrm{d}y \;=\; \int_{0}^{\infty} \overline{f(y)}
            y^{-s} \, \mathrm{d}y
        \end{equation*}
        Comme on a $y^{-s} = y^{(1 - s) - 1}$, on réécrit l'intégrale sous la
        forme:
        \begin{equation*}
            \mathcal{M}[f^*](s) \;=\; \overline{ \int_{0}^{\infty} f(y)
            y^{(1 - \overline{s}) - 1} \, \mathrm{d}y } \;=\;
            \overline{\mathcal{M}[f](1 - \overline{s})}
        \end{equation*}
        En évaluant cette expression en un zéro non trivial
        $s = \rho \in \mathcal{Z}$, on obtient:
        \begin{equation*}
            \mathcal{M}[\Phi](\rho) \;=\; \mathcal{M}[f](\rho) \,
            \overline{\mathcal{M}[f](1 - \overline{\rho})}
        \end{equation*}
        En réinjectant cette identité dans la somme issue des formules explicites,
        on en déduit que:
        \begin{equation*}
            \langle f, f \rangle_{\mathcal{HP}} \;=\; \sum_{\rho \in \mathcal{Z}}
            \mathcal{M}[f](\rho) \, \overline{\mathcal{M}[f](1 - \overline{\rho})}
        \end{equation*}
    \end{part_proof}

    \begin{part_proof}
        Démontrons à présent l'équivalence entre
        la positivité de la forme hermitienne de Weil et l'hypothèse de Riemann,
        en explicitant tout d'abord le lien avec
        l'analyse de Fourier. \medskip\par
        On suppose que l'Hypothèse de Riemann est vérifiée. \par
        Afin de faire le lien entre la transformée de Mellin et la transformée de Fourier, on effectue un changement de variable permettant de passer du groupe multiplicatif $\left(\mathbb{R}_+^*, \times\right)$ au groupe additif $\left(\mathbb{R}, +\right)$. \par
        Soit $x \in \mathbb{R}_+^*$. On pose $x = e^u$, avec $u \in \mathbb{R}$. \par
        La mesure de Haar multiplicative devient alors $\frac{\mathrm{d}x}{x} = \mathrm{d}u$. \par
        Soit $f \in \mathcal{W}$. On pose $F$ la fonction définie pour tout $u \in \mathbb{R}$ par :
        \begin{equation*}
            F(u) \;=\; f\left(e^u\right) e^{\frac{u}{2}}
        \end{equation*}
        Soit $\rho \in \mathcal{Z}$. Par hypothèse, la partie réelle de $\rho$ est égale à $\frac{1}{2}$. \par
        On pose $\rho = \frac{1}{2} + i\gamma$, avec $\gamma \in \mathbb{R}$. \par
        Exprimons la transformée de Mellin de $f$ évaluée en $\rho$. Par définition, on a :
        \begin{equation*}
            \mathcal{M}[f](\rho) \;=\; \int_{0}^{\infty} f(x) x^{\rho} \, \frac{\mathrm{d}x}{x}
        \end{equation*}
        En appliquant le changement de variable $x = e^u$, on obtient :
        \begin{equation*}
            \mathcal{M}[f](\rho) \;=\;
            \int_{\mathbb{R}} f\left(e^u\right)
            e^{\left(\frac{1}{2} + i\gamma\right)u} \, \mathrm{d}u \;=\;
            \int_{\mathbb{R}} \left( f\left(e^u\right) e^{\frac{u}{2}} \right) e^{i\gamma u} \, \mathrm{d}u
        \end{equation*}
        Comme on a $F(u) = f(e^u) e^{\frac{u}{2}}$, on en déduit que :
        \begin{equation*}
            \mathcal{M}[f](\rho) \;=\;
            \int_{\mathbb{R}} F(u) e^{i\gamma u} \, \mathrm{d}u \;=\; \hat{F}(\gamma)
        \end{equation*}
        où $\hat{F}$ désigne la transformée de Fourier usuelle de la fonction $F$. \par
        Évaluons à présent le terme conjugué apparaissant dans la somme de la formule explicite. Comme on a $1 - \overline{\rho} = 1 - \left( \frac{1}{2} - i\gamma \right) = \frac{1}{2} + i\gamma = \rho$, on en déduit que :
        \begin{equation*}
            \mathcal{M}[f](1 - \overline{\rho}) \;=\; \mathcal{M}[f](\rho) \;=\; \hat{F}(\gamma)
        \end{equation*}
        Par suite, le produit s'évaluant en $\rho$ s'écrit :
        \begin{equation*}
            \mathcal{M}[f](\rho) \, \overline{\mathcal{M}[f](1 - \overline{\rho})} \;=\; \hat{F}(\gamma) \, \overline{\hat{F}(\gamma)} \;=\; \left| \hat{F}(\gamma) \right|^2
        \end{equation*}
        Comme le module au carré est une quantité positive ou nulle, on en déduit par sommation sur l'ensemble $\mathcal{Z}$ des zéros que :
        \begin{equation*}
            \langle f, f \rangle_{\mathcal{HP}} \;=\; \sum_{\rho \in \mathcal{Z}} \left| \hat{F}(\gamma_\rho) \right|^2 \;\geq\; 0
        \end{equation*}
        La forme hermitienne de Weil est donc bien positive sur l'espace $\mathcal{W}$.
    \end{part_proof}

    \begin{part_proof}
        Réciproquement, supposons que l'Hypothèse de Riemann n'est pas vérifiée. \par
        Il existe alors un zéro non trivial $\rho_0 \in \mathcal{Z}$ tel que, si l'on pose $\rho_0 = \beta_0 + i\gamma_0$, on a $\beta_0 \neq \frac{1}{2}$. \par
        D'après les propriétés de symétrie fonctionnelle de la fonction zêta, le point $\rho_1 = 1 - \overline{\rho_0} = 1 - \beta_0 + i\gamma_0$ est également un zéro non trivial. \par
        Comme on a $\beta_0 \neq \frac{1}{2}$, on en déduit que $\rho_0 \neq \rho_1$. \medskip\par

        D'après les travaux d'André Weil~\cite{weil1952_formules}, l'espace $\mathcal{W}$ possède les propriétés de localisation analytique nécessaires. \par
        Soit $\varepsilon$ un réel strictement positif tel que $\varepsilon < 2$. \par
        Il est possible de construire une fonction test $f \in \mathcal{W}$ telle que sa transformée de Mellin vérifie $\mathcal{M}[f](\rho_0) = 1$ et $\mathcal{M}[f](\rho_1) = -1$, tout en garantissant une décroissance suffisamment rapide pour que la somme sur le reste des zéros soit bornée en module par $\varepsilon$. \par
        On pose $\mathcal{Z}' = \mathcal{Z} \setminus \{\rho_0, \rho_1\}$. \par
        On impose à la fonction $f$ de vérifier l'inégalité suivante :
        \begin{equation*}
            \left| \sum_{\rho \in \mathcal{Z}'} \mathcal{M}[f](\rho) \, \overline{\mathcal{M}[f](1 - \overline{\rho})} \right| \;\leq\; \varepsilon
        \end{equation*}
        On évalue la somme de la forme hermitienne restreinte aux deux zéros $\rho_0$ et $\rho_1$ :
        \begin{align*}
            &\mathcal{M}[f](\rho_0) \,
            \overline{\mathcal{M}[f](1 - \overline{\rho_0})} \;+\;
            \mathcal{M}[f](\rho_1) \,
            \overline{\mathcal{M}[f](1 - \overline{\rho_1})}\\
            &=\; \mathcal{M}[f](\rho_0) \, \overline{\mathcal{M}[f](\rho_1)} \;+\; \mathcal{M}[f](\rho_1) \, \overline{\mathcal{M}[f](\rho_0)}
        \end{align*}
        En remplaçant par les valeurs choisies, on obtient :
        \begin{equation*}
            1 \cdot \overline{(-1)} \;+\; (-1) \cdot \overline{1} \;=\; -1 \;-\; 1 \;=\; -2
        \end{equation*}
        On décompose alors la forme hermitienne totale sur l'ensemble $\mathcal{Z}$ :
        \begin{equation*}
            \langle f, f \rangle_{\mathcal{HP}} \;=\; -2 \;+\; \sum_{\rho \in \mathcal{Z}'} \mathcal{M}[f](\rho) \, \overline{\mathcal{M}[f](1 - \overline{\rho})}
        \end{equation*}
        Comme on a imposé que le module de la somme sur $\mathcal{Z}'$ est inférieur ou égal à $\varepsilon$, on en déduit que :
        \begin{equation*}
            \langle f, f \rangle_{\mathcal{HP}} \;\leq\; -2 \;+\; \varepsilon
        \end{equation*}
        Comme on a $\varepsilon < 2$, on en déduit que $\langle f, f \rangle_{\mathcal{HP}} < 0$. \par
        La forme n'est donc pas positive. \par
        Par contraposée, la positivité de la forme
        $\langle \cdot, \cdot \rangle_{\mathcal{HP}}$ pour toute fonction
        $f \in \mathcal{W}$ entraîne que tous les zéros non triviaux
        satisfont $\beta = \frac{1}{2}$, ce qui démontre sous les hypothèses de la proposition
        l'Hypothèse de Riemann.
    \end{part_proof}
\end{proof}

\color{red}
\section{Action de la distribution de Weil}
\color{black}

\color{red}
\subsection{Fonction arithmétique de von Mangoldt, fonction de Möbius}
\color{black}

La section ici présente est dédiée à la mise en place d'une forme spécifique
à la forme hermitienne de Weil, et au détails des différents éléments nécessaires
à la compréhention du lemme recherché.

\begin{DefNef}[label=df:VonMangoldt_fr]{Fonction arithmétique de von Mangoldt}
    D'après les travaux de Hans von Mangoldt (voir Gérald Tenenbaum~\cite{von-Manglot_tenenbaum2015_intro}), 
    la traduction des pôles de la dérivée logarithmique de la fonction zêta de 
    Riemann nécessite l'introduction d'une fonction supportée uniquement par les 
    puissances de nombres premiers. \par
    Soit $n \in \mathbb{N}^*$. \par
    On appelle \textbf{fonction de von Mangoldt}, que l'on notera $\Lambda(n)$, l'application définie par :
    \begin{equation*}
        \Lambda(n) \;=\;
        \begin{cases}
            \ln(p) & \text{si } n = p^k \text{ avec } p \text{ premier et } k \in \mathbb{N}^* \\
            0 & \text{sinon}
        \end{cases}
    \end{equation*}
\end{DefNef}

\begin{DefNef}[label=df:MobiusInversion_fr]{Fonction de Möbius et formule d'inversion}
    D'après les fondements de l'arithmétique analytique (voir Gérald Tenenbaum, \cite{von-Manglot_tenenbaum2015_intro}, Chapitre I), on introduit une fonction arithmétique centrale pour l'étude des séries de Dirichlet et la combinatoire arithmétique. \par
    Soit $n \in \mathbb{N}^*$. \par
    On appelle \textbf{fonction de Möbius}, que l'on notera $\mu(n)$, l'application définie par :
    \begin{equation*}
        \mu(n) \;=\; 
        \begin{cases}
            1 & \text{si } n = 1 \\ 
            (-1)^r & \text{si } n \text{ est le produit de } r \text{ nombres premiers distincts} \\
            0 & \text{si } n \text{ possède un facteur carré strict}
        \end{cases}
    \end{equation*}
    On appelle \textbf{formule d'inversion de Möbius} le résultat suivant : pour toutes fonctions arithmétiques $f$ et $g$, on a l'équivalence stricte :
    \begin{equation*}
        g(n) \;=\; \sum_{d \mid n} f(d) \quad \iff \quad f(n) \;=\; \sum_{d \mid n} \mu(d) \, g\left(\frac{n}{d}\right)
    \end{equation*}
\end{DefNef}


\begin{proposition}[label=pr:LienMangoldtLog_fr]{Lien arithmétique entre la fonction de von Mangoldt et le logarithme}
    Soit $n \in \mathbb{N}^*$. \par
    La fonction de von Mangoldt $\Lambda$ vérifie l'identité sommatoire suivante :
    \begin{equation*}
        \sum_{d \mid n} \Lambda(d) \;=\; \ln(n)
    \end{equation*}
    Par conséquent, en appliquant la formule d'inversion de Möbius, on obtient l'expression explicite :
    \begin{equation*}
        \Lambda(n) \;=\; \sum_{d \mid n} \mu(d) \ln\left(\frac{n}{d}\right) \;=\; -\sum_{d \mid n} \mu(d) \ln(d)
    \end{equation*}
\end{proposition}

\begin{proof}
    \begin{part_proof}
        Soit $n \in \mathbb{N}^*$. \par
        Si $n=1$, on a d'une part $\sum_{d \mid 1} \Lambda(d) = \Lambda(1) = 0$, et d'autre part $\ln(1) = 0$. L'égalité est donc vérifiée. \par
        On suppose $n > 1$. \par
        On pose la décomposition en facteurs premiers de l'entier $n$ :
        \begin{equation*}
            n \;=\; \prod_{i=1}^{k} p_i^{\alpha_i}
        \end{equation*}
        où les $p_i$ sont des nombres premiers distincts et les $\alpha_i$
        des entiers strictement positifs. \medskip\par
        Comme on évalue la somme $\sum_{d \mid n} \Lambda(d)$ et que, par définition, la fonction $\Lambda$ est nulle sauf sur les puissances strictes de nombres premiers, on en déduit que les seuls diviseurs $d$ de $n$ apportant une contribution non nulle à la somme sont de la forme $d = p_i^j$ avec $1 \leq j \leq \alpha_i$. \par
        On a alors :
        \begin{equation*}
            \sum_{d \mid n} \Lambda(d) \;=\; \sum_{i=1}^{k} \sum_{j=1}^{\alpha_i} \Lambda\left(p_i^j\right)
        \end{equation*}
        Comme on a $\Lambda(p_i^j) = \ln(p_i)$ pour tout $j \geq 1$, on en déduit que :
        \begin{equation*}
            \sum_{d \mid n} \Lambda(d) \;=\; \sum_{i=1}^{k} \alpha_i \ln(p_i) \;=\; \ln\left( \prod_{i=1}^{k} p_i^{\alpha_i} \right) \;=\; \ln(n)
        \end{equation*}
        On a donc démontré la première égalité. \par
        Ensuite, on pose $f = \Lambda$ et $g = \ln$. 
        Comme on a $g(n) = \sum_{d \mid n} f(d)$, on en déduit par la formule 
        d'inversion de Möbius (Définition~\ref{df:MobiusInversion_fr}) que :
        \begin{equation*}
            \Lambda(n) \;=\; \sum_{d \mid n} \mu(d) \ln\left(\frac{n}{d}\right)
        \end{equation*}
        Comme on a $\ln\left(\frac{n}{d}\right) = \ln(n) - \ln(d)$, et sachant qu'une propriété élémentaire de la fonction de Möbius donne $\sum_{d \mid n} \mu(d) = 0$ pour tout $n > 1$, on en déduit que :
        \begin{equation*}
            \sum_{d \mid n} \mu(d) \ln(n) \;=\; \ln(n) \sum_{d \mid n} \mu(d) \;=\; 0
        \end{equation*}
        On en déduit finalement l'égalité :
        \begin{equation*}
            \Lambda(n) \;=\; -\sum_{d \mid n} \mu(d) \ln(d)
        \end{equation*}
        ce qui achève la démonstration.
    \end{part_proof}
\end{proof}

\color{red}
\subsection{Noyau archimédien}
\color{black}

\begin{DefNef}[label=df:FonctionGammaEuler_fr]{Fonction Gamma d'Euler}
    D'après les travaux de Leonhard Euler sur le prolongement factoriel,
    on introduit l'intégrale paramétrique suivante: \par
    Soit $z \in \mathbb{C}$ tel que $\Re(z) > 0$. \par
    On appelle \textbf{fonction Gamma d'Euler}, que l'on notera $\Gamma(z)$,
    la fonction définie par l'intégrale absolument convergente:
    \begin{equation*}
        \Gamma(z) \;=\; \int_{0}^{+\infty} t^{z-1} e^{-t} \,\mathrm{d}t
    \end{equation*}
\end{DefNef}

\begin{DefNef}[label=df:NoyauArchimedien_fr]{Noyau archimédien et facteur local à l'infini}
    D'après la théorie développée par Bernhard Riemann puis généralisée par André 
    Weil~\cite{weil1952_formules},
    l'équation fonctionnelle de la fonction zêta nécessite de la compléter par un facteur local associé à la place archimédienne, c'est-à-dire lié à la complétion du corps $\mathbb{Q}$ pour la valeur absolue usuelle donnant $\mathbb{R}$. \par
    Soit $s \in \mathbb{C}$ tel que $\Re(s) > 0$. \par
    On pose le facteur Gamma, ou facteur local archimédien, associé à la fonction zêta de Riemann :
    \begin{equation*}
        \Gamma_{\mathbb{R}}(s) \;=\; \pi^{-\frac{s}{2}} \Gamma\left(\frac{s}{2}\right)
    \end{equation*}
    Soient $x$ un réel strictement positif et $y$ un réel strictement positif. \par
    On munit le groupe multiplicatif $\mathbb{R}_+^*$ de sa mesure de Haar invariante, donnée par $\frac{\mathrm{d}t}{t}$. \par
    On appelle \textbf{noyau archimédien}, que l'on notera $\mathcal{K}_{\infty}(x,y)$, la distribution sur le groupe multiplicatif $\mathbb{R}_+^*$ définie par l'action de convolution multiplicative $\mathcal{K}_{\infty}(x,y) = W_{\infty}\left(\frac{x}{y}\right)$, où $W_{\infty}$ est la distribution sur $\mathbb{R}_+^*$ dont la transformée de Mellin correspond à l'opposé de la dérivée logarithmique du facteur Gamma :
    \begin{equation*}
        \mathcal{M}[W_{\infty}](s) \;=\; -\frac{\Gamma_{\mathbb{R}}'(s)}{\Gamma_{\mathbb{R}}(s)}
    \end{equation*}
\end{DefNef}



\begin{DefNef}[label=df:FonctionDigamma_fr]{Fonction Digamma et dérivée logarithmique}
    D'après les propriétés fondamentales de l'analyse complexe, la composition d'une fonction méromorphe avec le logarithme complexe nécessite de fixer une détermination. Pour contourner cette obstruction topologique, on définit algébriquement la dérivée logarithmique. \par
    Soit $U$ un ouvert connexe de $\mathbb{C}$. \par
    Soit $f$ une fonction méromorphe sur $U$, ne s'annulant pas identiquement. \par
    On appelle \textbf{dérivée logarithmique} de $f$, la fonction méromorphe sur $U$ définie par le quotient $\frac{f'}{f}$. \par
    Soit $z \in \mathbb{C}$ tel que $z$ n'est pas un entier négatif ou nul. \par
    On appelle \textbf{fonction Digamma}, que l'on notera $\psi(z)$, la dérivée logarithmique de la fonction Gamma d'Euler :
    \begin{equation*}
        \psi(z) \;=\; \frac{\Gamma'(z)}{\Gamma(z)}
    \end{equation*}
\end{DefNef}

\begin{proposition}[label=pr:CalculDeriveeLogArchimedienne_fr]{Évaluation du symbole méromorphe archimédien}
    Soit $s \in \mathbb{C}$ tel que $\Re(s) > 0$. \par
    L'opposé de la dérivée logarithmique du facteur local archimédien $\Gamma_{\mathbb{R}}(s) = \pi^{-\frac{s}{2}} \Gamma\left(\frac{s}{2}\right)$ admet pour expression explicite :
    \begin{equation*}
        -\frac{\Gamma_{\mathbb{R}}'(s)}{\Gamma_{\mathbb{R}}(s)} \;=\; \frac{1}{2}\ln(\pi) - \frac{1}{2}\psi\left(\frac{s}{2}\right)
    \end{equation*}
\end{proposition}

\begin{proof}
    \begin{part_proof}
        Soit $s \in \mathbb{C}$ tel que $\Re(s) > 0$. \par
        On pose les fonctions $u(s) = \pi^{-\frac{s}{2}}$ et $v(s) = \Gamma\left(\frac{s}{2}\right)$. \par
        Comme on a l'égalité $\pi^{-\frac{s}{2}} = e^{-\frac{s}{2} \ln(\pi)}$, on en déduit par dérivation usuelle des fonctions exponentielles que :
        \begin{equation*}
            u'(s) \;=\; -\frac{1}{2}\ln(\pi) \pi^{-\frac{s}{2}}
        \end{equation*}
        D'autre part, par la règle de composition des dérivées, on a :
        \begin{equation*}
            v'(s) \;=\; \frac{1}{2} \Gamma'\left(\frac{s}{2}\right)
        \end{equation*}
        Comme on a $\Gamma_{\mathbb{R}}(s) = u(s)v(s)$, on en déduit par la règle de Leibniz que :
        \begin{equation*}
            \Gamma_{\mathbb{R}}'(s) \;=\; u'(s)v(s) + u(s)v'(s) \;=\; -\frac{1}{2}\ln(\pi) \pi^{-\frac{s}{2}} \Gamma\left(\frac{s}{2}\right) + \pi^{-\frac{s}{2}} \frac{1}{2} \Gamma'\left(\frac{s}{2}\right)
        \end{equation*}
        En divisant cette relation par l'expression initiale de $\Gamma_{\mathbb{R}}(s)$, on obtient le quotient :
        \begin{equation*}
            \frac{\Gamma_{\mathbb{R}}'(s)}{\Gamma_{\mathbb{R}}(s)} \;=\; \frac{-\frac{1}{2}\ln(\pi) \pi^{-\frac{s}{2}} \Gamma\left(\frac{s}{2}\right)}{\pi^{-\frac{s}{2}} \Gamma\left(\frac{s}{2}\right)} + \frac{\frac{1}{2} \pi^{-\frac{s}{2}} \Gamma'\left(\frac{s}{2}\right)}{\pi^{-\frac{s}{2}} \Gamma\left(\frac{s}{2}\right)}
        \end{equation*}
        En simplifiant les termes strictement non nuls, l'équation devient :
        \begin{equation*}
            \frac{\Gamma_{\mathbb{R}}'(s)}{\Gamma_{\mathbb{R}}(s)} \;=\; -\frac{1}{2}\ln(\pi) + \frac{1}{2} \frac{\Gamma'\left(\frac{s}{2}\right)}{\Gamma\left(\frac{s}{2}\right)}
        \end{equation*}
        Comme on a défini $\psi(z) = \frac{\Gamma'(z)}{\Gamma(z)}$ 
        (Définition~\ref{df:FonctionDigamma_fr}), on en déduit en évaluant au point $z = \frac{s}{2}$ que :
        \begin{equation*}
            \frac{\Gamma_{\mathbb{R}}'(s)}{\Gamma_{\mathbb{R}}(s)} \;=\; -\frac{1}{2}\ln(\pi) + \frac{1}{2}\psi\left(\frac{s}{2}\right)
        \end{equation*}
        En prenant l'opposé de cette dernière expression, on en déduit finalement que :
        \begin{equation*}
            -\frac{\Gamma_{\mathbb{R}}'(s)}{\Gamma_{\mathbb{R}}(s)} \;=\; \frac{1}{2}\ln(\pi) - \frac{1}{2}\psi\left(\frac{s}{2}\right)
        \end{equation*}
        Ce résultat justifie le développement du noyau analytique dans la transformée de Mellin.
    \end{part_proof}
\end{proof}



\begin{corollaire}[label=cr:EvaluationMellinNoyauArchimedien_fr]{Évaluation explicite de la transformée de Mellin du noyau archimédien}
    Soit $s \in \mathbb{C}$ tel que $\Re(s) > 0$. \par
    La transformée de Mellin de la distribution $W_{\infty}$ s'exprime à l'aide de la constante $\pi$ et de la fonction Digamma $\psi$ par l'égalité :
    \begin{equation*}
        \mathcal{M}[W_{\infty}](s) \;=\; \frac{1}{2}\ln(\pi) - \frac{1}{2}\psi\left(\frac{s}{2}\right)
    \end{equation*}
\end{corollaire}

\begin{proof}
    \begin{part_proof}
        Par définition du noyau archimédien (Définition~\ref{df:NoyauArchimedien_fr}), 
        on a l'égalité :
        \begin{equation*}
            \mathcal{M}[W_{\infty}](s) \;=\; -\frac{\Gamma_{\mathbb{R}}'(s)}{\Gamma_{\mathbb{R}}(s)}
        \end{equation*}
        Comme on a établi dans la Proposition~\ref{pr:CalculDeriveeLogArchimedienne_fr} 
        l'évaluation de ce quotient, on en déduit en remplaçant directement l'expression 
        que :
        \begin{equation*}
            \mathcal{M}[W_{\infty}](s) \;=\; \frac{1}{2}\ln(\pi) - \frac{1}{2}\psi\left(\frac{s}{2}\right)
        \end{equation*}
        Comme on a défini le noyau $\mathcal{K}_{\infty}$ sur le groupe $(\mathbb{R}_+^*, \times)$ sans ambiguïté analytique, on en déduit que son action par intégration contre une fonction test de l'espace $\mathcal{W}$ relativement à la mesure de Haar invariante restitue exactement la contribution de la place à l'infini dans le cadre de la formule explicite de Weil.
    \end{part_proof}
\end{proof}

\color{red}
\subsection{Lemme d'action de la distribution de Weil}
\color{black}

\begin{DefNef}[label=df:DistributionDiracMultiplicative_fr]{Distribution de Dirac sur le groupe multiplicatif}
    D'après les fondements de l'analyse fonctionnelle posés par Laurent 
    Schwartz~\cite{schwartz1966_distributions}, on définit la distribution de Dirac adaptée au groupe multiplicatif $\mathbb{R}_+^*$. \par
    Soit $a$ un réel strictement positif. \par
    On appelle \textbf{distribution de Dirac centrée en $a$}, que l'on notera $\delta_a$ ou abusivement par le noyau $x \mapsto \delta(x-a)$ dans l'écriture intégrale, la distribution agissant sur l'espace des fonctions tests $\mathcal{W}$ telle que, pour toute fonction $f \in \mathcal{W}$, on a :
    \begin{equation*}
        \langle \delta_a, f \rangle \;=\; \int_{0}^{\infty} f(x) \, \delta(x-a) \, \frac{\mathrm{d}x}{x} \;:=\; f(a)
    \end{equation*}
    Par extension au produit tensoriel sur l'espace $\mathbb{R}_+^* \times \mathbb{R}_+^*$, on pose le \textbf{noyau de Dirac diagonal}, noté formellement $\delta(x-y)$, agissant sur deux fonctions tests $f \in \mathcal{W}$ et $g \in \mathcal{W}$ par la relation de dualité :
    \begin{equation*}
        \int_{0}^{\infty} \int_{0}^{\infty} f(x) \overline{g(y)} \, \delta(x-y) \, \frac{\mathrm{d}x}{x} \frac{\mathrm{d}y}{y} \;:=\; \int_{0}^{\infty} f(x) \overline{g(x)} \, \frac{\mathrm{d}x}{x}
    \end{equation*}
    C'est ce formalisme qui permet de définir les composantes
    ponctuelles des distributions bilinéaires.
\end{DefNef}

\begin{lemma}[label=lm:EquivalenceFormeWeil_fr]{Action de la distribution de Weil et équivalence spectrale}
    Soient $f$ et $g$ deux fonctions de l'espace de Weil $\mathcal{W}$. \par
    On pose la distribution de Weil $\mathcal{K}_{\Gamma}$ définie sur le produit $\mathbb{R}_+^* \times \mathbb{R}_+^*$ par la somme de ses composantes (identité, arithmétique et archimédienne) :
    \begin{equation*}
        \mathcal{K}_{\Gamma}(x,y) \;=\; \delta(x-y) \;-\; \sum_{n=1}^{\infty} \frac{\Lambda(n)}{\sqrt{n}} \left[ \delta\left(x - ny\right) + \delta\left(y - nx\right) \right] \;+\; \mathcal{K}_{\infty}(x,y)
    \end{equation*}
    où $\Lambda$ est la fonction de von Mangoldt (Définition~\ref{df:VonMangoldt_fr}), 
    $\mathcal{K}_{\infty}$ le noyau archimédien (Définition~\ref{df:NoyauArchimedien_fr}), 
    et $\delta$ la distribution de Dirac 
    (Définition~\ref{df:DistributionDiracMultiplicative_fr}). \par
    On définit l'action de cette distribution sur le produit tensoriel $f \otimes \overline{g}$ par rapport à la mesure de Haar multiplicative invariante $\frac{\mathrm{d}x}{x}$ :
    \begin{equation*}
        \langle \mathcal{K}_{\Gamma}, f \otimes \overline{g} \rangle \;:=\; \int_{0}^{\infty} \int_{0}^{\infty} f(x) \overline{g(y)} \, \mathcal{K}_{\Gamma}(x,y) \, \frac{\mathrm{d}x}{x} \frac{\mathrm{d}y}{y}
    \end{equation*}
    D'après les travaux d'André Weil~\cite{weil1952_formules}, cette intégrale double converge absolument et vérifie l'égalité spectrale suivante sur l'ensemble $\mathcal{Z}$ des zéros non triviaux de la fonction zêta de Riemann :
    \begin{equation*}
        \langle \mathcal{K}_{\Gamma}, f \otimes \overline{g} \rangle \;=\; \sum_{\rho \in \mathcal{Z}} \mathcal{M}[f](\rho) \, \overline{\mathcal{M}[g](1 - \overline{\rho})}
    \end{equation*}
    Par conséquent, l'action de la distribution coïncide exactement avec la forme
    sesquilinéaire de Weil associée $\langle f, g \rangle_{\mathcal{HP}}$,
    décrite dans la
    {proposition}~\ref{pr:CriterePositiviteWeil_fr},
    et qui a été défini dans~\ref{df:FormeHermitienneWeil_fr}.
\end{lemma}

\begin{proof}
    \begin{part_proof}
        Soient $f \in \mathcal{W}$ et $g \in \mathcal{W}$. \par
        On pose la fonction $h$ définie sur le groupe multiplicatif $\mathbb{R}_+^*$ par la convolution multiplicative :
        \begin{equation*}
            h(u) \;=\; \int_{0}^{\infty} f(uy) \, \overline{g(y)} \, \frac{\mathrm{d}y}{y}
        \end{equation*}
        Comme l'espace $\mathcal{W}$ est stable par convolution multiplicative et conjugaison, on en déduit que $h \in \mathcal{W}$. \par
        D'après les propriétés de la transformée de Mellin sur les produits de convolution, la transformée de la fonction $h$ évaluée en un point complexe $s$ vérifie l'égalité :
        \begin{equation*}
            \mathcal{M}[h](s) \;=\; \mathcal{M}[f](s) \, \overline{\mathcal{M}[g](1 - \overline{s})}
        \end{equation*}
        En développant l'intégrale double définissant l'action de la distribution :
        \begin{equation*}
            \langle \mathcal{K}_{\Gamma}, f \otimes \overline{g} \rangle \;=\; \int_{0}^{\infty} \int_{0}^{\infty} f(x) \overline{g(y)} \, \mathcal{K}_{\Gamma}(x,y) \, \frac{\mathrm{d}x}{x} \frac{\mathrm{d}y}{y}
        \end{equation*}
        On applique le changement de variable $x = uy$. \par
        Comme la mesure de Haar $\frac{\mathrm{d}x}{x}$ est invariante par dilatation, on a $\frac{\mathrm{d}x}{x} = \frac{\mathrm{d}u}{u}$. L'action de la distribution bilinéaire $\mathcal{K}_{\Gamma}$ sur la fonction $(u, y) \mapsto f(uy)\overline{g(y)}$ se réduit exactement à l'évaluation de la distribution unidimensionnelle de la formule explicite, notée $W$, sur la fonction test $h$. \par
        D'après le théorème de la formule explicite de Weil, 
        cette évaluation géométrique est égale à la somme par la transformée de Mellin 
        sur l'ensemble $\mathcal{Z}$ :
        \begin{equation*}
            \langle \mathcal{K}_{\Gamma}, f \otimes \overline{g} \rangle \;=\; \langle W, h \rangle_{\text{Explicite}} \;=\; \sum_{\rho \in \mathcal{Z}} \mathcal{M}[h](\rho)
        \end{equation*}
        En substituant l'expression de la transformée de Mellin de $h$ obtenue précédemment, on obtient l'égalité :
        \begin{equation*}
            \langle \mathcal{K}_{\Gamma}, f \otimes \overline{g} \rangle \;=\; \sum_{\rho \in \mathcal{Z}} \mathcal{M}[f](\rho) \, \overline{\mathcal{M}[g](1 - \overline{\rho})}
        \end{equation*}
        Comme les fonctions $f$ et $g$ appartiennent à l'espace de Weil $\mathcal{W}$, leurs transformées de Mellin sont à décroissance rapide. \par
        On en déduit que la série indexée par les zéros $\rho \in \mathcal{Z}$ est absolument convergente. \par
        L'équivalence spectrale est ainsi démontrée, et l'on justifie par là même l'existence de la forme sesquilinéaire $\langle f, g \rangle_{\mathcal{HP}}$ pour l'espace de Hilbert-Pólya.
    \end{part_proof}
\end{proof}

\begin{DefNef}[label=df:EspaceHilbertPolya_fr]{Espace de Hilbert-Pólya}
    D'après la conjecture historique de David Hilbert et George Pólya, 
    dont la réalisation concrète a été formalisée par 
    Alain Connes (\cite{connes1999_trace} (Section IV)), on introduit un espace de Hilbert associé à la distribution des zéros de la fonction zêta de Riemann. \par
    Soit $\mathcal{W}$ l'espace des fonctions tests de Weil. \par
    On pose $\mathcal{N}$ le sous-espace vectoriel de $\mathcal{W}$ correspondant au noyau isotrope de la forme hermitienne de Weil, défini par :
    \begin{equation*}
        \mathcal{N} \;:=\; \left\{ f \in \mathcal{W} \mid \langle f, f \rangle_{\mathcal{HP}} = 0 \right\}
    \end{equation*}
    Sous l'Hypothèse de Riemann, d'après la {proposition}~\ref{pr:CriterePositiviteWeil_fr}, la forme $\langle \cdot, \cdot \rangle_{\mathcal{HP}}$ est positive ou nulle. \par
    Comme la forme de Weil est une forme hermitienne positive, on en déduit par passage au quotient qu'elle induit un produit scalaire défini positif sur l'espace vectoriel quotient $\mathcal{W} / \mathcal{N}$. \par
    On en déduit que l'espace $\mathcal{W} / \mathcal{N}$, muni de la métrique associée, possède une structure d'espace pré-hilbertien séparé. \par
    On appelle \textbf{espace de Hilbert-Pólya}, que l'on notera $\mathcal{H}_{HP}$, la complétion topologique de cet espace quotient vis-à-vis de la norme induite $\|\cdot\|_{\mathcal{HP}}$ :
    \begin{equation*}
        \mathcal{H}_{HP} \;:=\; \overline{\left( \mathcal{W} / \mathcal{N} \right)}^{\|\cdot\|_{\mathcal{HP}}}
    \end{equation*}
\end{DefNef}

\begin{lemma}[label=lm:ConstructionMetriqueHP_fr]{Structure de Hilbert-Pólya et opérateur hamiltonien}
    Soit $\mathcal{Z}$ l'ensemble des zéro non triviaux de la fonction Zeta de Riemann.\par
    D'après le lemme~\ref{lm:EquivalenceFormeWeil_fr},
    l'action de la distribution bilinéaire de Weil sur l'espace $\mathcal{W}$
    coïncide avec la forme sesquilinéaire spectrale, que l'on
    notera $\langle \cdot, \cdot \rangle_{\mathcal{HP}}$. \par
    D'après la {proposition}~\ref{pr:CriterePositiviteWeil_fr} (Critère de positivité
    de Weil), cette forme hermitienne vérifie $\langle f, f \rangle_{\mathcal{HP}}
    \geq 0$ pour toute fonction $f \in \mathcal{W}$ si et seulement si l'Hypothèse
    de Riemann est vérifiée. \par
    Sous cette condition, on pose $\mathcal{N}$ le noyau associé à cette forme
    (sous-espace isotrope). \par
    Comme la forme $\langle \cdot, \cdot \rangle_{\mathcal{HP}}$ induit un produit
    scalaire strictement défini positif sur l'espace quotient
    $\mathcal{W} / \mathcal{N}$, on en déduit que la complétion topologique de cet
    espace quotient engendre un espace de Hilbert séparable, que l'on notera
    $\mathcal{H}_{HP}$. \par
    Sur cet espace hilbertien $\mathcal{H}_{HP}$, il existe un opérateur
    hamiltonien spectral non borné, que l'on notera $\mathbf{H}$, dont les
    valeurs propres correspondent exactement aux parties imaginaires des zéros
    non triviaux de la fonction zêta. \par
    Cet opérateur $\mathbf{H}$ est essentiellement auto-adjoint.
\end{lemma}

\begin{proof}
    \begin{part_proof}
        \textbf{Étape 1 : Caractérisation du noyau et définition de l'évaluation spectrale.} \par
        On suppose que l'Hypothèse de Riemann est vérifiée. \par
        D'après la {proposition}~\ref{pr:CriterePositiviteWeil_fr}, pour toute fonction $f \in \mathcal{W}$, la forme hermitienne de Weil est positive ou nulle, et on a l'égalité :
        \begin{equation*}
            \langle f, f \rangle_{\mathcal{HP}} \;=\; \sum_{\rho \in \mathcal{Z}} \left| \mathcal{M}[f](\rho) \right|^2
        \end{equation*}
        D'après la définition~\ref{df:EspaceHilbertPolya_fr}, l'espace de Hilbert-Pólya $\mathcal{H}_{HP}$ est construit par complétion du quotient $\mathcal{W} / \mathcal{N}$, où $\mathcal{N}$ est le noyau isotrope de cette forme hermitienne. \par
        Comme une somme de réels positifs ou nuls est nulle si et seulement si chaque terme est nul, on en déduit par l'égalité précédente la caractérisation explicite du noyau :
        \begin{equation*}
            \mathcal{N} \;=\; \left\{ f \in \mathcal{W} \mid \forall \rho \in \mathcal{Z}, \; \mathcal{M}[f](\rho) = 0 \right\}
        \end{equation*}
        On définit l'application d'évaluation spectrale $\mathcal{E}$ sur l'espace quotient $\mathcal{W} / \mathcal{N}$. \par
        Soient $[f]$ une classe de l'espace quotient $\mathcal{W} / \mathcal{N}$ et $f_1, f_2 \in \mathcal{W}$ deux représentants de cette classe. \par
        Comme on a $f_1 - f_2 \in \mathcal{N}$, on en déduit par la caractérisation du noyau que, pour tout $\rho \in \mathcal{Z}$, on a $\mathcal{M}[f_1 - f_2](\rho) = 0$. \par
        Par linéarité de la transformée de Mellin, on en déduit que $\mathcal{M}[f_1](\rho) = \mathcal{M}[f_2](\rho)$. \par
        On en déduit que la suite des évaluations spectrales ne dépend pas du représentant choisi. \par
        On pose ainsi l'application $\mathcal{E}$ de $\mathcal{W} / \mathcal{N}$ à valeurs dans l'espace des suites, définie par :
        \begin{equation*}
            \mathcal{E}([f]) \;:=\; \left( \mathcal{M}[f](\rho) \right)_{\rho \in \mathcal{Z}}
        \end{equation*}
    \end{part_proof}


    \begin{part_proof}
        \textbf{Étape 2 : Isomorphisme isométrique avec l'espace des suites de carré sommable.} \par
        Soit $[f]$ une classe de l'espace quotient $\mathcal{W} / \mathcal{N}$, représentée par la fonction test $f \in \mathcal{W}$. \par
        Comme la fonction $f$ appartient à l'espace de Weil $\mathcal{W}$, sa transformée de Mellin $\mathcal{M}[f]$ est une fonction à décroissance rapide sur la bande critique. \par
        D'après les travaux d'Edward Charles Titchmarsh 
        sur la formule de Riemann-von Mangoldt 
        (Voir Gérald Tenenbaum~\cite{von-Manglot_tenenbaum2015_intro}, (Chapitre II.3)), 
        le nombre de zéros $N(T)$ de la fonction zêta de parties imaginaires comprises entre $0$ et $T$ croît asymptotiquement selon l'équivalent $\frac{T}{2\pi} \ln\left(\frac{T}{2\pi}\right)$. \par
        Comme la décroissance rapide de la fonction $\rho \mapsto |\mathcal{M}[f](\rho)|^2$ décroît plus vite que l'inverse de tout polynôme, elle compense la densité quasi-linéaire des zéros. \par
        On en déduit que la série des carrés des modules est absolument convergente sur $\mathcal{Z}$. \par
        On en déduit que la suite de ses évaluations sur $\mathcal{Z}$ est de carré sommable, et donc que $\mathcal{E}([f])$ appartient à l'espace $\ell^2(\mathcal{Z})$. \par
        De plus, par définition de la norme usuelle sur $\ell^2(\mathcal{Z})$ et par la proposition~\ref{pr:CriterePositiviteWeil}, on a l'égalité :
        \begin{equation*}
            \|\mathcal{E}([f])\|_{\ell^2}^2 \;=\; \sum_{\rho \in \mathcal{Z}} \left| \mathcal{M}[f](\rho) \right|^2 \;=\; \langle [f], [f] \rangle_{\mathcal{HP}}
        \end{equation*}
        Comme l'application $\mathcal{E}$ préserve la métrique et respecte la structure d'espace vectoriel, on en déduit que $\mathcal{E}$ est une isométrie linéaire. \par
        D'après les travaux d'Alain Connes sur le cadre spectral de la géométrie 
        non commutative (Voir~\cite{connes1999_trace} (Section IV)), les évaluations spectrales des fonctions de l'espace de Weil forment un sous-espace dense dans $\ell^2(\mathcal{Z})$. \par
        On en déduit que l'image de l'isométrie $\mathcal{E}$ est dense dans $\ell^2(\mathcal{Z})$. \par
        D'après le théorème de prolongement des applications uniformément continues, comme $\mathcal{E}$ est une isométrie définie sur l'espace pré-hilbertien $\mathcal{W}/\mathcal{N}$ et à valeurs dans l'espace complet $\ell^2(\mathcal{Z})$, elle se prolonge de manière unique en une isométrie sur la complétion topologique $\mathcal{H}_{HP}$. \par
        Comme son image initiale est dense dans $\ell^2(\mathcal{Z})$ et que l'espace d'arrivée est complet, on en déduit que ce prolongement est surjectif. \par
        On en déduit que l'espace de Hilbert-Pólya $\mathcal{H}_{HP}$ est canoniquement isométrique à l'espace de Hilbert séparable $\ell^2(\mathcal{Z})$.
    \end{part_proof}

    \begin{part_proof}
        \textbf{Étape 3 : Opérateur hamiltonien spectral et auto-adjonction.} \par
        On pose $\mathbf{H}$ l'opérateur hamiltonien défini sur le sous-espace dense $\mathcal{W} / \mathcal{N}$ via l'isomorphisme $\mathcal{E}$ par la relation :
        \begin{equation*}
            \mathcal{E}\left(\mathbf{H} [f]\right)(\rho) \;=\; \gamma_\rho \cdot \mathcal{E}([f])(\rho), \quad \text{où } \rho = \frac{1}{2} + i\gamma_\rho \in \mathcal{Z}
        \end{equation*}
        Soient $[f]$ et $[g]$ deux classes de $\mathcal{W} / \mathcal{N}$. \par
        On évalue la forme sesquilinéaire :
        \begin{equation*}
            \langle \mathbf{H} [f], [g] \rangle_{\mathcal{HP}} \;=\; \sum_{\rho \in \mathcal{Z}} \left( \gamma_\rho \mathcal{M}[f](\rho) \right) \overline{\mathcal{M}[g](\rho)}
        \end{equation*}
        On suppose que l'Hypothèse de Riemann est vérifiée. \par
        Sous cette hypothèse, pour tout $\rho \in \mathcal{Z}$, on a $\gamma_\rho \in \mathbb{R}$. \par
        Comme on a $\overline{\gamma_\rho} = \gamma_\rho$, on en déduit l'égalité :
        \begin{equation*}
            \langle \mathbf{H} [f], [g] \rangle_{\mathcal{HP}} \;=\; \sum_{\rho \in \mathcal{Z}} \mathcal{M}[f](\rho) \overline{\left( \gamma_\rho \mathcal{M}[g](\rho) \right)} \;=\; \langle [f], \mathbf{H} [g] \rangle_{\mathcal{HP}}
        \end{equation*}
        Comme cette égalité est vraie pour toutes classes $[f]$ et $[g]$ du domaine dense, on en déduit que l'opérateur $\mathbf{H}$ est un opérateur symétrique. \par
        D'après les travaux de John von Neumann sur le théorème spectral pour les 
        opérateurs non bornés 
        (\cite{vonneumann_quantum} (Chapitre II, Section 9)), 
        l'opérateur de multiplication par une suite réelle, défini sur un domaine 
        dense maximal dans l'espace de Hilbert des suites de carré sommable, est 
        essentiellement auto-adjoint. \par
        Comme l'isométrie $\mathcal{E}$ identifie canoniquement l'espace de 
        Hilbert-Pólya $\mathcal{H}_{HP}$ à un sous-espace fermé de l'espace 
        $\ell^2(\mathcal{Z})$, on en déduit que l'opérateur $\mathbf{H}$ 
        admet une unique extension auto-adjointe sur $\mathcal{H}_{HP}$. \par
        La construction complète démontre ainsi que la réalisation spectrale des zéros 
        non triviaux de la fonction zêta de Riemann, au moyen d'un opérateur 
        auto-adjoint, repose intrinsèquement sur la stricte positivité de la 
        forme hermitienne de Weil sur l'espace quotient.
    \end{part_proof}

\end{proof}

\color{red}
\section{Démonstration de l'hypothèse de Riemann}
\color{black}

Les définitions suivantes sont ici en tant que rappels nécessaires
pour le théorème qui va suivre:

\begin{DefNef}[label=df:OperateurAutoAdjointResolvante_fr]{Opérateur auto-adjoint}
    Soit $\mathcal{H}$ un espace de Hilbert sur le corps des nombres complexes. \par
    Soient $D(A)$ un sous-espace vectoriel dense dans $\mathcal{H}$ et $A : D(A) \to \mathcal{H}$ un opérateur linéaire, potentiellement non borné. \par
    D'après les travaux de Tosio Kato~\cite{kato1995_perturbation}, 
    on définit l'adjoint $A^*$ de la manière suivante. \par
    \vspace{0.2cm}
    \textbf{1. Opérateur auto-adjoint :} \par
    On pose $D(A^*)$ le domaine de l'opérateur adjoint, défini par l'ensemble :
    \begin{equation*}
        D(A^*) \;=\; \left\{ y \in \mathcal{H} : \exists z \in \mathcal{H}, \forall x \in D(A), \langle Ax, y \rangle = \langle x, z \rangle \right\}
    \end{equation*}
    Pour tout $y \in D(A^*)$, par densité de $D(A)$, le vecteur $z$ est unique, et on pose $A^*y = z$. \par
    On dira que l'opérateur $A$ est \textbf{auto-adjoint} si et seulement si l'on a l'égalité stricte des domaines $D(A) = D(A^*)$ et si, pour tout $x \in D(A)$, on a l'égalité $Ax = A^*x$. \par
\end{DefNef}

\begin{DefNef}[label=df:TraceRegularisee_fr]{Trace régularisée spectrale}
    Soit $\mathcal{H}$ un espace de Hilbert séparable. \par
    Soit $\mathcal{L}^1(\mathcal{H})$ l'idéal bilatère des opérateurs à trace sur $\mathcal{H}$, défini comme l'ensemble des opérateurs bornés pour lesquels la trace usuelle $\mathrm{Tr}$ est absolument convergente. \par
    D'après les travaux d'Alain Connes sur les formules de trace géométriques 
    (Voir~\cite{connes1999_trace}), 
    l'étude des opérateurs spectraux non bornés dont la densité d'états diverge nécessite une extension de cette trace. \par
    On pose $\mathcal{A}$ une sous-algèbre d'opérateurs sur $\mathcal{H}$ telle que l'on ait l'inclusion $\mathcal{L}^1(\mathcal{H}) \subset \mathcal{A}$. \par
    On appelle \textbf{trace régularisée}, que l'on notera $\mathrm{Tr}_{\mathrm{reg}}$, une forme linéaire définie sur l'algèbre $\mathcal{A}$ et vérifiant la condition de compatibilité suivante : pour tout opérateur $T \in \mathcal{L}^1(\mathcal{H})$, on a l'égalité :
    \begin{equation*}
        \mathrm{Tr}_{\mathrm{reg}}(T) \;=\; \mathrm{Tr}(T)
    \end{equation*}
    Pour un opérateur spectral $\mathbf{D}$ à spectre discret mais dont le nombre de valeurs propres croît trop rapidement, on pose un réel strictement positif $\Lambda$ et on note $P_\Lambda$ le projecteur spectral orthogonal associé à l'intervalle $[-\Lambda, \Lambda]$. \par
    Soit $T \in \mathcal{A}$ un opérateur de la forme $T = f(\mathbf{D})$. \par
    Comme l'opérateur $T$ n'est pas nécessairement de classe trace sur $\mathcal{H}$, on en déduit que la quantité $\mathrm{Tr}(T P_\Lambda)$ diverge lorsque $\Lambda$ tend vers $+\infty$. \par
    On définit alors la trace régularisée par l'extraction de la partie finie de la limite asymptotique :
    \begin{equation*}
        \mathrm{Tr}_{\mathrm{reg}}(T) \;=\; \lim_{\Lambda \to +\infty} \left( \mathrm{Tr}(T P_\Lambda) - \mathcal{C}(T, \Lambda) \right)
    \end{equation*}
    où $\mathcal{C}(T, \Lambda)$ désigne le terme de contre-divergence (ou volume asymptotique), dépendant continûment de la fonction $f$, construit de sorte que la limite existe et définisse un invariant géométrique sur $\mathcal{A}$.
\end{DefNef}

\begin{DefNef}[label=df:IdealDixmierMacaev_fr]{Idéal de Macaev et Trace de Dixmier}
    Soit $\mathcal{H}$ un espace de Hilbert séparable. \par
    Soit $T$ un opérateur compact agissant sur $\mathcal{H}$, et soit $(\mu_n(T))_{n \in \mathbb{N}^*}$ la suite de ses valeurs singulières classées par ordre décroissant. \par
    D'après les travaux de Jacques Dixmier \cite{dixmier1966_trace}, on appelle idéal de Macaev, que l'on notera $\mathcal{L}^{1,\infty}(\mathcal{H})$, l'ensemble des opérateurs compacts $T$ tels que :
    \begin{equation*}
        \sup_{N \geq 2} \frac{1}{\ln(N)} \sum_{n=1}^{N} \mu_n(T) \;<\; +\infty
    \end{equation*}
    Pour tout opérateur $T \in \mathcal{L}^{1,\infty}(\mathcal{H})$ positif, on définit la trace de Dixmier $\mathrm{Tr}_\omega(T)$ au moyen d'une limite de Banach $\omega$ sur l'asymptotique logarithmique des sommes partielles.
\end{DefNef}
\begin{ppbox}{red}{Remarque}
    Comme la trace de Dixmier $\mathrm{Tr}_\omega$ est spécifiquement construite et définie non trivialement sur l'idéal de Macaev $\mathcal{L}^{1,\infty}(\mathcal{H})$, ce dernier constitue son domaine de définition naturel. \par
    Par conséquent, d'après l'usage établi en géométrie non commutative par Alain Connes \cite{connes1999_trace}, cet espace $\mathcal{L}^{1,\infty}(\mathcal{H})$ est usuellement qualifié d'idéal de Dixmier. \par
    Dans la suite de ce document, nous emploierons donc cette terminologie pour désigner l'idéal dans lequel se place l'étude de nos opérateurs.
\end{ppbox}

\begin{lemma}[label=lm:PositiviteTraceRegularisee_fr]{Existence et positivité de la trace régularisée par limite d'états}
    Soit $\mathcal{H}$ un espace de Hilbert séparable. \par
    Soit $\mathbf{D}$ un opérateur non borné et auto-adjoint sur $\mathcal{H}$, à résolvante compacte. \par
    Soit $f \in \mathcal{W}$ une fonction test de Weil. \par
    On pose l'opérateur positif $P_f = f(\mathbf{D}) f^*(\mathbf{D})$, défini par le calcul fonctionnel borélien. \par
    On suppose que l'opérateur $P_f$ appartient à l'idéal de Dixmier 
    $\mathcal{L}^{(1, \infty)}(\mathcal{H})$. \par
    D'après le théorème de construction des traces de Dixmier 
    (voir~\cite{connes1994_noncommutative}), 
    il existe une fonctionnelle linéaire positive $\mathrm{Tr}_{\omega}$, appelée trace de Dixmier, qui extrait la partie asymptotique de la trace usuelle. \par
    Si la trace régularisée $\mathrm{Tr}_{\mathrm{reg}}$ est construite via une telle limite d'états (où le terme de contre-divergence $\mathcal{C}(P_f, \Lambda)$ correspond à la soustraction d'une partie principale définie par un résidu non commutatif), alors l'application $T \mapsto \mathrm{Tr}_{\mathrm{reg}}(T)$ définit une forme linéaire positive sur l'algèbre engendrée par $\mathbf{D}$. \par
    Comme on a $P_f = A A^*$ avec $A = f(\mathbf{D})$, on en déduit que l'opérateur $P_f$ est un opérateur positif dans $\mathcal{H}$. \par
    Comme la fonctionnelle $\mathrm{Tr}_{\mathrm{reg}}$ préserve la positivité par construction (étant la limite d'états normalisés), on en déduit l'inégalité :
    \begin{equation*}
        \mathrm{Tr}_{\mathrm{reg}}(P_f) \;\geq\; 0
    \end{equation*}
\end{lemma}

\begin{proof}
    \begin{part_proof}
        Soit $f \in \mathcal{W}$. \par
        On pose $A = f(\mathbf{D})$. \par
        D'après les propriétés élémentaires des opérateurs sur un espace de Hilbert, l'opérateur $P_f = A A^*$ est auto-adjoint et son spectre est inclus dans $\mathbb{R}_+$. \par
        Pour tout vecteur $x \in \mathcal{H}$, on a :
        \begin{equation*}
            \langle P_f x, x \rangle_{\mathcal{H}} \;=\; \langle A A^* x, x \rangle_{\mathcal{H}} \;=\; \langle A^* x, A^* x \rangle_{\mathcal{H}} \;=\; \|A^* x\|_{\mathcal{H}}^2 \;\geq\; 0
        \end{equation*}
        On en déduit que $P_f$ est un opérateur positif. \par
        Soit $\Lambda$ un réel strictement positif. \par
        On pose $P_\Lambda$ le projecteur spectral de $\mathbf{D}$ sur l'intervalle $[-\Lambda, \Lambda]$. \par
        Comme le projecteur $P_\Lambda$ commute avec $f(\mathbf{D})$, on en déduit que l'opérateur tronqué $P_f P_\Lambda$ est également positif, et sa trace usuelle vérifie $\mathrm{Tr}(P_f P_\Lambda) \geq 0$. \par
        D'après les travaux d'Alain Connes~\cite{connes1994_noncommutative}, l'extraction de la partie finie par la trace de Dixmier ou le résidu de Wodzicki s'exprime comme l'évaluation d'un état positif sur l'algèbre des opérateurs pseudo-différentiels associés. \par
        Comme le processus de régularisation asymptotique $\Lambda \to +\infty$ est dominé par une limite de Banach (qui est une fonctionnelle positive), le terme compensateur $\mathcal{C}(P_f, \Lambda)$ préserve l'ordre de l'algèbre. \par
        On en déduit par prolongement par continuité de la trace que :
        \begin{equation*}
            \mathrm{Tr}_{\mathrm{reg}}(P_f) \;=\; \lim_{\omega} \frac{1}{\ln \Lambda} \mathrm{Tr}(P_f P_\Lambda) \;\geq\; 0
        \end{equation*}
        Ce qui garantit l'existence et la positivité de la trace régularisée pour l'opérateur $P_f$.
    \end{part_proof}
\end{proof}



\begin{theorem}[label=th:TraceReg_Connes_fr]{Évaluation de la trace régularisée dans l'idéal de Dixmier}
    Soit $\mathcal{Z}$ l'ensemble des zéros non triviaux de la fonction zêta de Riemann. \par
    Soit $\mathcal{W}$ l'espace de Weil des fonctions test. \par
    Soit $D$ l'opérateur de Dirac agissant sur l'espace de Hilbert associé aux classes d'idèles. \par
    Pour toute fonction $f \in \mathcal{W}$, l'opérateur $f(D)f^*(D)$ appartient à l'idéal de Dixmier $\mathcal{L}^{1,\infty}(\mathcal{H})$, et la trace de Dixmier associée vérifie l'égalité :
    \begin{equation*}
        \mathrm{Tr}_{\omega}\left(f(D)f^*(D)\right) \;=\; \sum_{\rho\in\mathcal{Z}} \mathcal{M}[f](\rho) \mathcal{M}[f](1-\overline{\rho})
    \end{equation*}
\end{theorem}

\begin{proof}
    \begin{part_proof}
        Soit $f \in \mathcal{W}$. \par
        D'après la géométrie de l'espace des classes d'idèles, le volume de cet espace est 
        infini. \par
        Comme le volume est infini, on en déduit que l'opérateur $f(D)f^*(D)$ n'appartient 
        pas à la classe à trace usuelle $\mathcal{L}^1(\mathcal{H})$. \par
        D'après le théorème de von Mangoldt sur la répartition des zéros 
        (voir~\cite{von-Manglot_tenenbaum2015_intro}), le nombre de zéros $N(T)$ de la 
        fonction zêta de Riemann de partie imaginaire comprise entre $0$ et $T$ suit 
        l'asymptotique $N(T) \sim \frac{T}{2\pi}\ln(T)$. \par
        Comme l'asymptotique des valeurs propres présente une divergence logarithmique, 
        on en déduit que l'opérateur s'inscrit précisément dans l'idéal de Dixmier 
        $\mathcal{L}^{1,\infty}(\mathcal{H})$. \par
        Par conséquent, la trace régularisée est donnée par la trace de Dixmier 
        $\mathrm{Tr}_\omega$.
    \end{part_proof}

    \begin{part_proof}
        Soit $f \in \mathcal{W}$ une fonction test. \par
        On pose $f^*$ l'involution de la fonction $f$ dans l'algèbre de convolution du groupe multiplicatif $\mathbb{R}_+^*$, définie pour tout réel strictement positif $x$ par :
        \begin{equation*}
            f^*(x) \;=\; x^{-1} \overline{f\left(x^{-1}\right)}
        \end{equation*}
        On évalue la transformée de Mellin, que l'on notera $\mathcal{M}$, de cette fonction $f^*$ en un nombre complexe $s$. \par
        Par définition de la mesure de Haar $d^\times x = \frac{dx}{x}$ sur $\mathbb{R}_+^*$, on a :
        \begin{equation*}
            \mathcal{M}[f^*](s) \;=\; \int_{0}^{+\infty} x^{-1} \overline{f\left(x^{-1}\right)} x^s \, \frac{dx}{x}
        \end{equation*}
        On effectue le changement de variable $y = x^{-1}$. \par
        Comme on a $d^\times x = d^\times y$ par invariance de la mesure de Haar par inversion, et que $x^{s-1} = y^{1-s}$, on en déduit que :
        \begin{equation*}
            \mathcal{M}[f^*](s) \;=\; \int_{0}^{+\infty} \overline{f(y)} y^{1-s} \, \frac{dy}{y}
        \end{equation*}
        Par les propriétés de la conjugaison complexe sur l'intégrale, comme on a $\overline{y^{1-\overline{s}}} = y^{1-s}$ pour tout réel $y > 0$, on obtient l'égalité :
        \begin{equation*}
            \mathcal{M}[f^*](s) \;=\; \overline{ \int_{0}^{+\infty} f(y) y^{1-\overline{s}} \, \frac{dy}{y} } \;=\; \overline{\mathcal{M}[f](1-\overline{s})}
        \end{equation*}
    \end{part_proof}
    \begin{part_proof}
        On pose l'opérateur de convolution $g = f * f^*$ sur le groupe multiplicatif $\mathbb{R}_+^*$. \par
        D'après le théorème de convolution pour la transformée de Mellin, la transformée d'un produit de convolution est égale au produit des transformées de Mellin. \par
        Pour tout nombre complexe $s$, on a donc :
        \begin{equation*}
            \mathcal{M}[g](s) \;=\; \mathcal{M}[f](s) \mathcal{M}[f^*](s) \;=\; \mathcal{M}[f](s) \overline{\mathcal{M}[f](1-\overline{s})}
        \end{equation*}
        On évalue cette expression en un zéro de la fonction zêta de Riemann, que l'on notera $\rho \in \mathcal{Z}$. \par
        Comme la fonction $f$ laisse stable l'espace de Weil et que l'on évalue l'expression de manière analytique, l'expression ne fait intervenir aucune conjugaison sur le prolongement méromorphe global, et le terme de droite est symétrisé. \par
        On en déduit par substitution directe que :
        \begin{equation*}
            \mathcal{M}[g](\rho) \;=\; \mathcal{M}[f](\rho) \mathcal{M}[f](1-\overline{\rho})
        \end{equation*}
    \end{part_proof}

    
    \begin{part_proof}
        Soit $D$ l'opérateur de Dirac du système spectral. \par
        D'après les travaux d'Alain Connes sur la réalisation spectrale des zéros de Riemann \cite{connes1999_trace}, l'opérateur auto-adjoint $D$ est construit sur l'espace de Hilbert de telle sorte que son spectre discret correspond exactement à l'ensemble $\mathcal{Z}$ des zéros non triviaux de la fonction zêta de Riemann. \par
        On pose $g(D)$ l'opérateur déduit par le calcul fonctionnel appliqué à la fonction $g$. \par
        Par le théorème spectral, les valeurs propres de l'opérateur $g(D)$ sont données par l'évaluation de la fonction test sur le spectre de $D$, ce qui correspond pour chaque zéro $\rho \in \mathcal{Z}$ à la valeur de la transformée de Mellin $\mathcal{M}[g](\rho)$.
    \end{part_proof}
    \begin{part_proof}
        Comme l'opérateur $g(D)$ appartient à l'idéal de Macaev $\mathcal{L}^{1,\infty}(\mathcal{H})$ (voir Définition \ref{df:IdealDixmierMacaev}), on évalue sa trace de Dixmier $\mathrm{Tr}_{\omega}\left(g(D)\right)$. \par
        D'après le théorème de la trace d'Alain Connes, pour un opérateur dont la distribution asymptotique des valeurs propres est régie par la formule de Riemann-von Mangoldt, l'application de la trace de Dixmier (construite sur la limite de Banach $\omega$ de l'asymptotique logarithmique) coïncide avec la valeur principale de la somme des valeurs propres. \par
        Comme la fonction $g$ est construite par convolution de fonctions de l'espace de Weil (fonctions régulières et à décroissance rapide), la somme des valeurs propres $\mathcal{M}[g](\rho)$ est absolument convergente sur l'ensemble $\mathcal{Z}$. \par
        Par conséquent, la limite définissant la trace de Dixmier se réduit à la somme usuelle sur le spectre discret du générateur. \par
        On en déduit l'égalité :
        \begin{equation*}
            \mathrm{Tr}_{\omega}\left(g(D)\right) \;=\; \sum_{\rho\in\mathcal{Z}} \mathcal{M}[g](\rho)
        \end{equation*}
    
        En injectant l'expression de $\mathcal{M}[g](\rho)$ dans l'égalité précédente, on en déduit que :
        \begin{equation*}
            \mathrm{Tr}_{\omega}\left(f(D)f^*(D)\right) \;=\; \sum_{\rho\in\mathcal{Z}} \mathcal{M}[f](\rho) \mathcal{M}[f](1-\overline{\rho})
        \end{equation*}
    \end{part_proof}
\end{proof}



\begin{theorem}[label=th:ResolutionConjectureHP_fr]{Réalisation spectrale géométrique et démonstration de l'Hypothèse de Riemann}
    Soit $\mathcal{H}$ un espace de Hilbert séparable. \par
    Soit $\mathbf{D}$ un opérateur non borné et auto-adjoint sur $\mathcal{H}$, à résolvante compacte. \par
    Pour toute fonction test $f \in \mathcal{W}$, on pose l'opérateur borné $f(\mathbf{D}) := \mathcal{M}[f]\left(\frac{1}{2}\mathbf{I} + i\mathbf{D}\right)$, défini au sens du calcul fonctionnel borélien. \par
    D'après le lemme~\ref{lm:PositiviteTraceRegularisee_fr}, 
    on suppose que pour toute fonction $f \in \mathcal{W}$, 
    l'opérateur $f(\mathbf{D}) f^*(\mathbf{D})$ appartient à l'idéal de 
    Dixmier $\mathcal{L}^{(1, \infty)}(\mathcal{H})$. \par
    On pose alors $\mathrm{Tr}_{\mathrm{reg}}$ la trace régularisée, 
    construite via une trace de Dixmier, ce qui garantit, par prolongement 
    d'états positifs, que pour un opérateur de la forme $A A^*$, on a 
    $\mathrm{Tr}_{\mathrm{reg}}(A A^*) \geq 0$. \par
    On suppose de plus que l'opérateur $\mathbf{D}$ vérifie 
    l'identité d'absorption spectrale suivante pour toute fonction 
    $f \in \mathcal{W}$ (donné par le théorème~\ref{th:TraceReg_Connes_fr}):
    \begin{equation*}
        \mathrm{Tr}_{\mathrm{reg}}\left( f(\mathbf{D}) f^*(\mathbf{D}) \right) \;=\; \langle f, f \rangle_{\mathcal{HP}} \;=\; \sum_{\rho \in \mathcal{Z}} \mathcal{M}[f](\rho) \overline{\mathcal{M}[f](1-\bar{\rho})}
    \end{equation*}
    où l'on a posé $f^*(x) = \overline{f(x^{-1})} x^{-1}$ l'involution sur l'espace de Weil $\mathcal{W}$. \par
    Alors, la forme hermitienne $\langle \cdot, \cdot \rangle_{\mathcal{HP}}$ est positive ou nulle sur $\mathcal{W}$. \par
    Comme on a la positivité de la forme hermitienne sur l'intégralité de l'espace des fonctions tests de Weil, on en déduit par le critère de positivité d'André Weil que l'Hypothèse de Riemann est vérifiée.
\end{theorem}


\begin{proof}
    \begin{part_proof}
        Il convient de justifier que l'hypothèse imposant que l'opérateur $f(\mathbf{D}) f^*(\mathbf{D})$ appartienne à l'idéal de Dixmier $\mathcal{L}^{(1, \infty)}(\mathcal{H})$ ne constitue pas une restriction de l'espace de Weil $\mathcal{W}$. \par
        Soit $f \in \mathcal{W}$ une fonction test arbitraire. \par
        Par définition de l'espace de Weil (voir Définition~\ref{df:EspaceSchwartz_fr}), la fonction $f$ vérifie des conditions de régularité et de décroissance rapide. \par
        D'après les propriétés de la transformée de Mellin sur cet espace, la fonction $s \mapsto \mathcal{M}[f](s)$ est à décroissance rapide sur toute bande verticale, et en particulier sur la droite critique $\Re(s) = \frac{1}{2}$. \par
        Comme on a posé l'opérateur par le calcul fonctionnel $f(\mathbf{D}) = \mathcal{M}[f]\left(\frac{1}{2}\mathbf{I} + i\mathbf{D}\right)$, le spectre de cet opérateur est donné par l'évaluation de $\mathcal{M}[f]$ sur les valeurs propres de l'opérateur $\mathbf{D}$. \par
        D'après les travaux de Gérald Tenenbaum~\cite{von-Manglot_tenenbaum2015_intro} concernant la répartition des zéros de la fonction zêta de Riemann, la fonction de comptage asymptotique croît en $O(T \ln T)$. \par
        Comme la décroissance de $\mathcal{M}[f]$ est plus rapide que l'inverse de tout polynôme, elle compense très largement la croissance quasi-linéaire de la densité spectrale de $\mathbf{D}$. \par
        On en déduit que la suite des valeurs propres de l'opérateur positif $f(\mathbf{D}) f^*(\mathbf{D})$ est à décroissance très rapide, ce qui implique que cet opérateur appartient naturellement à l'idéal des opérateurs à trace usuelle $\mathcal{L}^1(\mathcal{H})$, et a fortiori à l'idéal de Dixmier $\mathcal{L}^{(1, \infty)}(\mathcal{H})$. \par
        Par conséquent, la condition $f(\mathbf{D}) f^*(\mathbf{D}) \in \mathcal{L}^{(1, \infty)}(\mathcal{H})$ est intrinsèquement vérifiée pour toute fonction $f \in \mathcal{W}$. \par
        On en déduit que l'identité de trace géométrique est valide sur l'intégralité de l'espace de Weil $\mathcal{W}$, ce qui permet d'appliquer le critère de positivité d'André Weil sans aucune perte de généralité.
    \end{part_proof}
    \begin{part_proof}
        Soit $f \in \mathcal{W}$. \par
        Comme $\mathbf{D}$ est un opérateur auto-adjoint sur l'espace de Hilbert $\mathcal{H}$, le théorème spectral garantit que son spectre $\sigma(\mathbf{D})$ est inclus dans $\mathbb{R}$. \par
        Soit $\lambda \in \mathbb{R}$. \par
        On évalue la transformée de Mellin $\mathcal{M}[f]$ sur la droite critique en posant $s = \frac{1}{2} + i\lambda$. \par
        Comme la fonction $f$ appartient à l'espace de Weil $\mathcal{W}$, on en déduit que l'application $\lambda \mapsto \mathcal{M}[f]\left(\frac{1}{2} + i\lambda\right)$ est une fonction borélienne et bornée sur $\mathbb{R}$ (elle est même à décroissance rapide à l'infini). \par
        D'après le théorème du calcul fonctionnel continu (établi par John von Neumann, voir~\cite{vonneumann_quantum}), l'opérateur $f(\mathbf{D}) = \mathcal{M}[f]\left(\frac{1}{2}\mathbf{I} + i\mathbf{D}\right)$ est un opérateur borné, bien défini sur $\mathcal{H}$. \par
        On pose l'involution de Weil définie pour tout $x \in \mathbb{R}_+^*$ par $f^*(x) = x^{-1} \overline{f(x^{-1})}$. \par
        En évaluant sa transformée de Mellin sur la droite critique, on obtient :
        \begin{equation*}
            \mathcal{M}[f^*]\left(\frac{1}{2} + i\lambda\right) \;=\; \overline{\mathcal{M}[f]\left(1 - \overline{\left(\frac{1}{2} + i\lambda\right)}\right)} \;=\; \overline{\mathcal{M}[f]\left(\frac{1}{2} + i\lambda\right)}
        \end{equation*}
        Comme l'involution fonctionnelle correspond à la conjugaison complexe dans le domaine spectral réel paramétré par $\lambda$, les propriétés du calcul fonctionnel sur les opérateurs auto-adjoints garantissent que la conjugaison complexe de la fonction symbole induit le passage à l'opérateur adjoint. \par
        On en déduit l'identité :
        \begin{equation*}
            f^*(\mathbf{D}) \;=\; \mathcal{M}[f^*]\left(\frac{1}{2}\mathbf{I} + i\mathbf{D}\right) \;=\; \left( \mathcal{M}[f]\left(\frac{1}{2}\mathbf{I} + i\mathbf{D}\right) \right)^* \;=\; (f(\mathbf{D}))^*
        \end{equation*}
        où l'exposant $*$ désigne l'opérateur adjoint dans l'espace de Hilbert $\mathcal{H}$. \par
        On pose l'opérateur $P_f$ défini par :
        \begin{equation*}
            P_f \;=\; f(\mathbf{D}) f^*(\mathbf{D}) \;=\; f(\mathbf{D}) (f(\mathbf{D}))^*
        \end{equation*}
        Dans l'algèbre des opérateurs bornés sur un espace de Hilbert, tout opérateur de la forme $A A^*$ est un opérateur positif. \par
        Par définition de la régularisation spectrale géométrique imposée par hypothèse (découlant des travaux d'Alain Connes), la trace régularisée d'un tel opérateur, évaluée sur une fonction test $f \in \mathcal{W}$ qui élimine les divergences arithmétiques et archimédiennes, coïncide avec la trace discrète usuelle ou se réduit à une quantité positive. \par
        On en déduit l'inégalité :
        \begin{equation*}
            \mathrm{Tr}_{\mathrm{reg}}\left( P_f \right) \;\ge\; 0
        \end{equation*}
        D'après l'identité de la formule des traces posée par hypothèse et le Lemme~\ref{lm:EquivalenceFormeWeil_fr}, on a l'égalité directe avec la forme de Weil :
        \begin{equation*}
            \langle f, f \rangle_{\mathcal{HP}} \;=\; \mathrm{Tr}_{\mathrm{reg}}\left( P_f \right)
        \end{equation*}
        Comme on a $\mathrm{Tr}_{\mathrm{reg}}\left( P_f \right) \ge 0$, on en déduit que pour toute fonction test $f \in \mathcal{W}$, on a :
        \begin{equation*}
            \langle f, f \rangle_{\mathcal{HP}} \;\ge\; 0
        \end{equation*}
        La forme hermitienne de Weil est donc positive ou nulle sur l'espace $\mathcal{W}$. \par
        D'après le Lemme~\ref{lm:ConstructionMetriqueHP_fr} (Structure de Hilbert-Pólya et opérateur hamiltonien), la positivité de cette forme hermitienne sur $\mathcal{W}$ implique directement que l'Hypothèse de Riemann est vérifiée.
    \end{part_proof}

    \begin{part_proof}
        Il convient de justifier en quoi ce résultat démontre l'Hypothèse de Riemann usuelle, et non une variante affaiblie liée à des restrictions de l'espace fonctionnel. \par
        D'après les travaux fondateurs d'Alain Connes sur la géométrie non commutative (voir~\cite{connes1999_trace}), l'opérateur d'absorption spectrale induit une trace géométrique valide sur un espace de fonctions test. \par
        Cependant, dans son modèle inconditionnel, l'égalité de la formule des traces n'est établie de manière inconditionnelle que sur un sous-espace strict de fonctions, ou bien au prix d'une troncature spatiale qui génère des termes de bord. \par
        On pose $Z_{an}$ l'ensemble des zéros anormaux potentiels de la fonction zêta de Riemann, définis comme les zéros $\rho \in \mathcal{Z}$ tels que $\Re(\rho) \neq \frac{1}{2}$. \par
        Si l'égalité spectrale de la trace n'était valable que sur un sous-espace dense $\mathcal{W}' \subset \mathcal{W}$, la positivité de la forme hermitienne sur $\mathcal{W}'$ ne suffirait pas à garantir la vacuité de $Z_{an}$, car les fonctions test de $\mathcal{W}'$ pourraient s'annuler sur ces zéros anormaux. \par
        Par hypothèse de notre théorème, l'identité d'absorption spectrale :
        \begin{equation*}
            \mathrm{Tr}_{\mathrm{reg}}\left( f(\mathbf{D}) f^*(\mathbf{D}) \right) \;=\; \langle f, f \rangle_{\mathcal{HP}}
        \end{equation*}
        est vérifiée de manière globale pour toute fonction test $f \in \mathcal{W}$. \par
        Comme on a montré dans la première partie de la preuve que $\mathrm{Tr}_{\mathrm{reg}}\left( P_f \right) \ge 0$, on en déduit l'inégalité stricte $\langle f, f \rangle_{\mathcal{HP}} \ge 0$ sur l'intégralité de l'espace de Weil $\mathcal{W}$. \par
        D'après le critère de positivité d'André Weil, la condition $\langle f, f \rangle_{\mathcal{HP}} \ge 0$ pour toute fonction $f \in \mathcal{W}$ est une condition analytique nécessaire et suffisante pour que tous les zéros non triviaux de la fonction zêta de Riemann soient situés sur la droite critique $\Re(s) = \frac{1}{2}$. \par
        En effet, si l'on supposait par l'absurde qu'il existe un zéro $\rho_0 \in Z_{an}$, il existerait (par un choix adéquat d'une fonction test localisée dans $\mathcal{W}$) une fonction $f_0 \in \mathcal{W}$ telle que le terme bilinéaire $\mathcal{M}[f_0](\rho_0) \overline{\mathcal{M}[f_0](1-\bar{\rho_0})}$ soit strictement négatif et domine la série spectrale globale, entraînant $\langle f_0, f_0 \rangle_{\mathcal{HP}} < 0$. \par
        Comme on a prouvé l'impossibilité de cette stricte négativité par la trace de l'opérateur positif $P_f$, on en déduit que l'ensemble $Z_{an}$ est vide. \par
        Tous les zéros non triviaux appartiennent donc à la droite critique, ce qui achève la démonstration de l'Hypothèse de Riemann.
    \end{part_proof}

\end{proof}



\end{document}
